\begin{titlepage}
    \centering
    \setstretch{1.0} % Title page usually uses single spacing for grouping

    % 1. Institution Info (12pt)
    {\fontsize{12pt}{14pt}\selectfont 
    TARTU ÜLIKOOL \\
    Sotsiaalteaduste Valdkond \\
    Narva Kolledž \par}

    \vspace{2cm}

    % 2. Curriculum (12pt)
    {\fontsize{12pt}{14pt}\selectfont 
    Infotehnoloogiliste süsteemide arendus\par}

    \vspace{2.5cm}

    % 3. Author (12pt)
    {\fontsize{12pt}{14pt}\selectfont Sergei Ivanov \par}

    \vspace{1cm}

    % 4. Title (16pt, BOLD, ALL CAPS)
    % Note: Subtitles should also follow the size/bold rules if part of the main title
    {\fontsize{16pt}{20pt}\selectfont \textbf{Neon Fractal Tree} \par}

    \vspace{1cm}

    % 5. Type of work (12pt)
    {\fontsize{12pt}{14pt}\selectfont Kodutöö Project Nr 1 \par}

    \vspace{3cm}

    % 6. Supervisor (12pt) - Position/Degree required
    \begin{flushright}
    {\fontsize{12pt}{14pt}\selectfont 
    Õppejõud: Jaak Vilo, \\
    Andmeteaduse magistriõppekava programmijuht\par}
    \end{flushright}

    \vfill

    % 7. Location and Year (12pt)
    {\fontsize{12pt}{14pt}\selectfont Narva 2026 \par}

\end{titlepage}

\begin{abstract}
See report käsitleb Project NR 1 - Neon Fractal Tree, mis on loodud Pythonis, kasutades turtle graafikateeki. Projekt keskendub rekursiivse algoritmi rakendamisele, et luua visuaalselt köitev fraktaalpuu, mille värvid muutuvad dünaamiliselt. Kood on optimeeritud nii, et see jookseb sujuvalt isegi suurema sügavusega, vältides liigset ressursikasutust. Lõpptulemus on esteetiliselt meeldiv ja näitab, kuidas programmeerimine võib olla kunstiline väljendusvorm.
\end{abstract}

\vspace{0.5cm}

\noindent\textbf{Märksõnad:} Python, turtle graphics, recursion, fractal, visualization

\tableofcontents
\clearpage

